
\documentclass[11pt]{article}

\usepackage[margin=1in]{geometry}
\usepackage{amsmath,amssymb,amsthm}
\usepackage{graphicx}
\usepackage{booktabs}
\usepackage{subcaption}
\usepackage[numbers]{natbib}
\usepackage{xcolor}
\usepackage{float}
\usepackage{longtable}

\title{A Reduced WNT--RA--HOX Model for Colorectal Cancer Stemness}
\author{Amir, Nate, PK}
\date{June 2026}

\begin{document}

\maketitle

\section{Model Overview}

This model describes the regulatory interactions among the WNT, HOX, and
retinoic acid (RA) pathways in colorectal cancer. The model contains seven
state variables representing key molecular components involved in stemness,
differentiation, and treatment response:
\[
B(t)=\text{$\beta$-catenin}, \qquad
P(t)=\text{APC}, \qquad
H_5(t)=\text{HOXA5},
\]
\[
H_{13}(t)=\text{HOXA13}, \qquad
M(t)=\text{MYC}, \qquad
R(t)=\text{RA}, \qquad
C(t)=\text{CYP26A1}.
\]

The model is built around three major biological feedback loops. The first is
a pro-stemness loop:
\[
\beta\text{-catenin}
\longrightarrow
\text{MYC}
\longrightarrow
\text{HOXA13}
\longrightarrow
\beta\text{-catenin}.
\]
The second is a differentiation loop:
\[
\text{RA}
\longrightarrow
\text{HOXA5}
\longrightarrow
\text{APC}
\dashv
\beta\text{-catenin}.
\]
The third is an RA homeostasis loop:
\[
\text{RA}
\longrightarrow
\text{CYP26A1}
\dashv
\text{RA}.
\]

Reduced APC functionality is used to represent APC-mutant colorectal cancer
cells. APC loss weakens the suppression of $\beta$-catenin, leading to
increased MYC and HOXA13 activity, reduced differentiation, and increased
stemness.

\section{Dimensional Model}

The dimensional WNT--RA--HOX--CYP model is
\begin{align}
\frac{dB}{dt}
&=
a_BW
+
v_{13}\frac{H_{13}^{n_H}}{K_{13}^{n_H}+H_{13}^{n_H}}
-
d_BB
-
k_{PB}PB
-
v_5\frac{H_5B}{K_5+B},
\\
\frac{dP}{dt}
&=
a_P
\frac{1+\rho_5H_5}
{1+\rho_BB+\rho_{13}H_{13}}
-
\Big(d_{P0}+d_{P1}(1-\Theta_P)\Big)P,
\\
\frac{dH_5}{dt}
&=
a_5
+
v_R\frac{R}{K_R+R}
-
d_5H_5
-
v_M\frac{MH_5}{K_M+M},
\\
\frac{dH_{13}}{dt}
&=
a_{13}
+
v_{B13}\frac{B^{n_B}}{K_{B13}^{n_B}+B^{n_B}}
+
v_{M13}\frac{M^{n_M}}{K_{M13}^{n_M}+M^{n_M}}
-
d_{13}H_{13},
\\
\frac{dM}{dt}
&=
a_M
+
v_{BM}\frac{B^{n_B}}{K_{BM}^{n_B}+B^{n_B}}
-
d_MM,
\\
\frac{dR}{dt}
&=
u_R(t)
-
d_RR
-
k_{CR}CR,
\\
\frac{dC}{dt}
&=
a_C
+
v_{RC}\frac{R}{K_{RC}+R}
+
v_{BC}\frac{B^{n_B}}{K_{BC}^{n_B}+B^{n_B}}
-
d_CC.
\end{align}

\[
u_R(t)
=
u_0
+
\frac{D}{2}
\left[
\tanh(q(t-t_1))
-
\tanh(q(t-t_2))
\right].
\]

\section{Nondimensional Model}

The resulting nondimensional system is
\[
b=\frac{B}{B_0},\quad
p=\frac{P}{P_0},\quad
h_5=\frac{H_5}{H_{50}},\quad
h_{13}=\frac{H_{13}}{H_{130}},
\]
\[
m=\frac{M}{M_0},\quad
r=\frac{R}{R_0},\quad
c=\frac{C}{C_0},\quad
\tau=d_Bt.
\]

\begin{align}
\frac{db}{d\tau}
&=
w
+
\eta_{13}\frac{h_{13}^{n_H}}{\kappa_{13}^{n_H}+h_{13}^{n_H}}
-
b
-
\lambda_Ppb
-
\lambda_5\frac{h_5b}{\kappa_5+b},
\\
\frac{dp}{d\tau}
&=
\frac{1}{\varepsilon_P}
\left[
\frac{1+\rho_5h_5}
{1+\rho_Bb+\rho_{13}h_{13}}
-
\delta_P(\theta_P)p
\right],
\\
\frac{dh_5}{d\tau}
&=
\frac{1}{\varepsilon_5}
\left[
a_5
+
\eta_R\frac{r}{\kappa_R+r}
-
h_5
-
\eta_M\frac{mh_5}{\kappa_M+m}
\right],
\\
\frac{dh_{13}}{d\tau}
&=
\frac{1}{\varepsilon_{13}}
\left[
a_{13}
+
\eta_{B13}\frac{b^{n_B}}{\kappa_{B13}^{n_B}+b^{n_B}}
+
\eta_{M13}\frac{m^{n_M}}{\kappa_{M13}^{n_M}+m^{n_M}}
-
h_{13}
\right],
\\
\frac{dm}{d\tau}
&=
\frac{1}{\varepsilon_M}
\left[
a_M
+
\eta_{BM}\frac{b^{n_B}}{\kappa_{BM}^{n_B}+b^{n_B}}
-
m
\right],
\\
\frac{dr}{d\tau}
&=
\frac{1}{\varepsilon_R}
\left[
\mu_R(\tau)
-
r
-
\lambda_Ccr
\right],
\\
\frac{dc}{d\tau}
&=
\frac{1}{\varepsilon_C}
\left[
a_C
+
\eta_{RC}\frac{r}{\kappa_{RC}+r}
+
\eta_{BC}\frac{b^{n_B}}{\kappa_{BC}^{n_B}+b^{n_B}}
-
c
\right].
\end{align}

\[
\delta_P(\theta_P)=1+\delta_{P1}(1-\theta_P).
\]

\[
S(\tau)
=
\frac{
b(\tau)\left[1+\alpha_{13}h_{13}(\tau)\right]
}
{
1+p(\tau)+\alpha_5h_5(\tau)
}.
\]
where $\theta_P$ measures functional APC activity.

The nondimensional treatment function is
\begin{equation}
\mu_R(\tau)
=
\mu_0
+
\frac{D_R}{2}
\left[
\tanh\!\big(q(\tau-\tau_1)\big)
-
\tanh\!\big(q(\tau-\tau_2)\big)
\right].
\end{equation}

\section{Expected Model Behavior}

The model is designed to reproduce four key biological behaviors:

\begin{enumerate}
\item \textbf{Normal tissue:} high APC, controlled $\beta$-catenin, moderate MYC and HOXA13, and low stemness.

\item \textbf{APC-mutant tissue:} reduced APC functionality increases $\beta$-catenin signaling, which elevates MYC and HOXA13 and suppresses differentiation.

\item \textbf{RA treatment:} increased RA induces HOXA5, enhances APC activity, suppresses $\beta$-catenin signaling, and lowers stemness.

\item \textbf{CYP-mediated resistance:} elevated CYP26A1 accelerates RA degradation, reducing the effectiveness of RA-based therapies.
\end{enumerate}

To quantify cellular phenotype, the stemness index is defined by
\begin{equation}
S(\tau)
=
\frac{
b(\tau)\left[1+\alpha_{13}h_{13}(\tau)\right]
}
{
1+p(\tau)+\alpha_5h_5(\tau)
}.
\end{equation}

Large values of $S$ correspond to highly proliferative, stem-like states,
whereas small values correspond to differentiated cellular states.



\newpage

\section{Dimensional and Nondimensional Model}

We consider the seven dimensional variables
\[
B(t),\quad P(t),\quad H_5(t),\quad H_{13}(t),\quad M(t),\quad R(t),\quad C(t),
\]
where \(B\) denotes \(\beta\)-catenin, \(P\) denotes APC, \(H_5\) denotes HOXA5, \(H_{13}\) denotes HOXA13, \(M\) denotes MYC, \(R\) denotes retinoic acid, and \(C\) denotes CYP26A1.

We introduce the dimensionless variables
\[
b=\frac{B}{B_0},\qquad
p=\frac{P}{P_0},\qquad
h_5=\frac{H_5}{H_{50}},\qquad
h_{13}=\frac{H_{13}}{H_{130}},
\]
\[
m=\frac{M}{M_0},\qquad
r=\frac{R}{R_0},\qquad
c=\frac{C}{C_0},
\]
and the dimensionless time
\[
\tau=d_B t.
\]

Thus,
\[
\frac{d}{dt}=d_B\frac{d}{d\tau}.
\]

\subsection{1. \(\beta\)-catenin equation}

The dimensional equation is
\[
\frac{dB}{dt}
=
a_BW
+
v_{13}
\frac{H_{13}^{n_H}}
{K_{13}^{n_H}+H_{13}^{n_H}}
-
d_BB
-
k_{PB}PB
-
v_5\frac{H_5B}{K_5+B}.
\]

Substitute
\[
B=B_0b,\qquad P=P_0p,\qquad H_5=H_{50}h_5,\qquad H_{13}=H_{130}h_{13}.
\]

Then
\[
d_BB_0\frac{db}{d\tau}
=
a_BW
+
v_{13}
\frac{H_{130}^{n_H}h_{13}^{n_H}}
{K_{13}^{n_H}+H_{130}^{n_H}h_{13}^{n_H}}
-
d_BB_0b
-
k_{PB}P_0B_0pb
-
v_5
\frac{H_{50}h_5B_0b}{K_5+B_0b}.
\]

Divide by \(d_BB_0\):
\[
\frac{db}{d\tau}
=
\frac{a_BW}{d_BB_0}
+
\frac{v_{13}}{d_BB_0}
\frac{h_{13}^{n_H}}
{\left(\frac{K_{13}}{H_{130}}\right)^{n_H}+h_{13}^{n_H}}
-
b
-
\frac{k_{PB}P_0}{d_B}pb
-
\frac{v_5H_{50}}{d_BB_0}
\frac{h_5b}{\frac{K_5}{B_0}+b}.
\]

Define
\[
w=\frac{a_BW}{d_BB_0},\qquad
\eta_{13}=\frac{v_{13}}{d_BB_0},\qquad
\kappa_{13}=\frac{K_{13}}{H_{130}},
\]
\[
\lambda_P=\frac{k_{PB}P_0}{d_B},\qquad
\lambda_5=\frac{v_5H_{50}}{d_BB_0},\qquad
\kappa_5=\frac{K_5}{B_0}.
\]

The nondimensional equation is
\[
\boxed{
\frac{db}{d\tau}
=
w
+
\eta_{13}
\frac{h_{13}^{n_H}}
{\kappa_{13}^{n_H}+h_{13}^{n_H}}
-
b
-
\lambda_Ppb
-
\lambda_5\frac{h_5b}{\kappa_5+b}.
}
\]

\subsection{2. APC equation}

The dimensional equation is
\[
\frac{dP}{dt}
=
a_P
\frac{1+\rho_5H_5}
{1+\rho_BB+\rho_{13}H_{13}}
-
\Big(\gamma_{P0}+\gamma_{P1}(1-\Theta_P)\Big)P.
\]

We use \(\gamma_{P0}\) and \(\gamma_{P1}\) for APC degradation rates to avoid confusion with the APC variable \(P\).

Substitute
\[
P=P_0p,\qquad B=B_0b,\qquad H_5=H_{50}h_5,\qquad H_{13}=H_{130}h_{13}.
\]

Then
\[
d_BP_0\frac{dp}{d\tau}
=
a_P
\frac{1+\rho_5H_{50}h_5}
{1+\rho_BB_0b+\rho_{13}H_{130}h_{13}}
-
\Big(\gamma_{P0}+\gamma_{P1}(1-\Theta_P)\Big)P_0p.
\]

Divide by \(d_BP_0\):
\[
\frac{dp}{d\tau}
=
\frac{a_P}{d_BP_0}
\frac{1+\rho_5H_{50}h_5}
{1+\rho_BB_0b+\rho_{13}H_{130}h_{13}}
-
\frac{\gamma_{P0}+\gamma_{P1}(1-\Theta_P)}{d_B}p.
\]

Choose
\[
P_0=\frac{a_P}{\gamma_{P0}},
\qquad
\varepsilon_P=\frac{d_B}{\gamma_{P0}}.
\]

Then
\[
\frac{a_P}{d_BP_0}
=
\frac{\gamma_{P0}}{d_B}
=
\frac{1}{\varepsilon_P}.
\]

Define
\[
\rho_5^*=\rho_5H_{50},\qquad
\rho_B^*=\rho_BB_0,\qquad
\rho_{13}^*=\rho_{13}H_{130},
\]
and, for simplicity, relabel
\[
\rho_5^*\to \rho_5,\qquad
\rho_B^*\to \rho_B,\qquad
\rho_{13}^*\to \rho_{13}.
\]

Also define
\[
\delta_{P1}=\frac{\gamma_{P1}}{\gamma_{P0}},
\qquad
\theta_P=\Theta_P,
\]
so that
\[
\delta_P(\theta_P)=1+\delta_{P1}(1-\theta_P).
\]

The nondimensional equation is
\[
\boxed{
\frac{dp}{d\tau}
=
\frac{1}{\varepsilon_P}
\left[
\frac{1+\rho_5h_5}
{1+\rho_Bb+\rho_{13}h_{13}}
-
\delta_P(\theta_P)p
\right].
}
\]

\subsection{3. HOXA5 equation}

The dimensional equation is
\[
\frac{dH_5}{dt}
=
a_5
+
v_R\frac{R}{K_R+R}
-
d_5H_5
-
v_M\frac{MH_5}{K_M+M}.
\]

Substitute
\[
H_5=H_{50}h_5,\qquad R=R_0r,\qquad M=M_0m.
\]

Then
\[
d_BH_{50}\frac{dh_5}{d\tau}
=
a_5
+
v_R\frac{R_0r}{K_R+R_0r}
-
d_5H_{50}h_5
-
v_M\frac{M_0mH_{50}h_5}{K_M+M_0m}.
\]

Divide by \(d_BH_{50}\):
\[
\frac{dh_5}{d\tau}
=
\frac{a_5}{d_BH_{50}}
+
\frac{v_R}{d_BH_{50}}
\frac{r}{\frac{K_R}{R_0}+r}
-
\frac{d_5}{d_B}h_5
-
\frac{v_M}{d_B}
\frac{mh_5}{\frac{K_M}{M_0}+m}.
\]

Choose
\[
H_{50}=\frac{a_5}{d_5},
\qquad
\varepsilon_5=\frac{d_B}{d_5}.
\]

Then
\[
\frac{a_5}{d_BH_{50}}=\frac{d_5}{d_B}=\frac{1}{\varepsilon_5}.
\]

Define
\[
\eta_R=\frac{v_R}{d_5H_{50}},
\qquad
\kappa_R=\frac{K_R}{R_0},
\qquad
\eta_M=\frac{v_M}{d_5},
\qquad
\kappa_M=\frac{K_M}{M_0}.
\]

The nondimensional equation is
\[
\boxed{
\frac{dh_5}{d\tau}
=
\frac{1}{\varepsilon_5}
\left[
1
+
\eta_R\frac{r}{\kappa_R+r}
-
h_5
-
\eta_M\frac{mh_5}{\kappa_M+m}
\right].
}
\]

If one does not choose \(H_{50}=a_5/d_5\), the first term becomes a free dimensionless basal production parameter \(a_5^*\), giving
\[
\frac{dh_5}{d\tau}
=
\frac{1}{\varepsilon_5}
\left[
a_5^*
+
\eta_R\frac{r}{\kappa_R+r}
-
h_5
-
\eta_M\frac{mh_5}{\kappa_M+m}
\right].
\]

\subsection{4. HOXA13 equation}

The dimensional equation is
\[
\frac{dH_{13}}{dt}
=
a_{13}
+
v_{B13}
\frac{B^{n_B}}
{K_{B13}^{n_B}+B^{n_B}}
+
v_{M13}
\frac{M^{n_M}}
{K_{M13}^{n_M}+M^{n_M}}
-
d_{13}H_{13}.
\]

Substitute
\[
H_{13}=H_{130}h_{13},\qquad B=B_0b,\qquad M=M_0m.
\]

Then
\[
d_BH_{130}\frac{dh_{13}}{d\tau}
=
a_{13}
+
v_{B13}
\frac{B_0^{n_B}b^{n_B}}
{K_{B13}^{n_B}+B_0^{n_B}b^{n_B}}
+
v_{M13}
\frac{M_0^{n_M}m^{n_M}}
{K_{M13}^{n_M}+M_0^{n_M}m^{n_M}}
-
d_{13}H_{130}h_{13}.
\]

Divide by \(d_BH_{130}\):
\[
\frac{dh_{13}}{d\tau}
=
\frac{a_{13}}{d_BH_{130}}
+
\frac{v_{B13}}{d_BH_{130}}
\frac{b^{n_B}}
{\left(\frac{K_{B13}}{B_0}\right)^{n_B}+b^{n_B}}
+
\frac{v_{M13}}{d_BH_{130}}
\frac{m^{n_M}}
{\left(\frac{K_{M13}}{M_0}\right)^{n_M}+m^{n_M}}
-
\frac{d_{13}}{d_B}h_{13}.
\]

Choose
\[
H_{130}=\frac{a_{13}}{d_{13}},
\qquad
\varepsilon_{13}=\frac{d_B}{d_{13}}.
\]

Define
\[
\eta_{B13}=\frac{v_{B13}}{d_{13}H_{130}},
\qquad
\kappa_{B13}=\frac{K_{B13}}{B_0},
\]
\[
\eta_{M13}=\frac{v_{M13}}{d_{13}H_{130}},
\qquad
\kappa_{M13}=\frac{K_{M13}}{M_0}.
\]

The nondimensional equation is
\[
\boxed{
\frac{dh_{13}}{d\tau}
=
\frac{1}{\varepsilon_{13}}
\left[
1
+
\eta_{B13}
\frac{b^{n_B}}
{\kappa_{B13}^{n_B}+b^{n_B}}
+
\eta_{M13}
\frac{m^{n_M}}
{\kappa_{M13}^{n_M}+m^{n_M}}
-
h_{13}
\right].
}
\]

If one does not choose \(H_{130}=a_{13}/d_{13}\), the basal term becomes \(a_{13}^*\).

\subsection{5. MYC equation}

The dimensional equation is
\[
\frac{dM}{dt}
=
a_M
+
v_{BM}
\frac{B^{n_B}}
{K_{BM}^{n_B}+B^{n_B}}
-
d_MM.
\]

Substitute
\[
M=M_0m,\qquad B=B_0b.
\]

Then
\[
d_BM_0\frac{dm}{d\tau}
=
a_M
+
v_{BM}
\frac{B_0^{n_B}b^{n_B}}
{K_{BM}^{n_B}+B_0^{n_B}b^{n_B}}
-
d_MM_0m.
\]

Divide by \(d_BM_0\):
\[
\frac{dm}{d\tau}
=
\frac{a_M}{d_BM_0}
+
\frac{v_{BM}}{d_BM_0}
\frac{b^{n_B}}
{\left(\frac{K_{BM}}{B_0}\right)^{n_B}+b^{n_B}}
-
\frac{d_M}{d_B}m.
\]

Choose
\[
M_0=\frac{a_M}{d_M},
\qquad
\varepsilon_M=\frac{d_B}{d_M}.
\]

Define
\[
\eta_{BM}=\frac{v_{BM}}{d_MM_0},
\qquad
\kappa_{BM}=\frac{K_{BM}}{B_0}.
\]

The nondimensional equation is
\[
\boxed{
\frac{dm}{d\tau}
=
\frac{1}{\varepsilon_M}
\left[
1
+
\eta_{BM}
\frac{b^{n_B}}
{\kappa_{BM}^{n_B}+b^{n_B}}
-
m
\right].
}
\]

If one does not choose \(M_0=a_M/d_M\), the basal term becomes \(a_M^*\).

\subsection{6. Retinoic acid equation}

The dimensional equation is
\[
\frac{dR}{dt}
=
u_R(t)
-
d_RR
-
k_{CR}CR.
\]

The input is
\[
u_R(t)
=
u_0
+
A_d
\left[
1+
\cos\left(
\frac{2\pi t}{T_d}-\phi
\right)
\right]
+
\frac{D}{2}
\left[
\tanh(q(t-t_1))
-
\tanh(q(t-t_2))
\right].
\]

Substitute
\[
R=R_0r,\qquad C=C_0c,\qquad \tau=d_Bt.
\]

Then
\[
d_BR_0\frac{dr}{d\tau}
=
u_R(t)
-
d_RR_0r
-
k_{CR}C_0cR_0r.
\]

Divide by \(d_BR_0\):
\[
\frac{dr}{d\tau}
=
\frac{u_R(t)}{d_BR_0}
-
\frac{d_R}{d_B}r
-
\frac{k_{CR}C_0}{d_B}cr.
\]

Define
\[
\varepsilon_R=\frac{d_B}{d_R},
\qquad
\lambda_C=\frac{k_{CR}C_0}{d_R},
\qquad
\mu_R(\tau)=\frac{u_R(t)}{d_RR_0}.
\]

Then
\[
\frac{u_R(t)}{d_BR_0}
=
\frac{1}{\varepsilon_R}\mu_R(\tau).
\]

The nondimensional equation is
\[
\boxed{
\frac{dr}{d\tau}
=
\frac{1}{\varepsilon_R}
\left[
\mu_R(\tau)
-
r
-
\lambda_Ccr
\right].
}
\]

Now define the nondimensional input parameters
\[
\mu_0=\frac{u_0}{d_RR_0},
\qquad
A_R=\frac{A_d}{d_RR_0},
\qquad
D_R=\frac{D}{d_RR_0},
\]
\[
T_R=d_BT_d,
\qquad
q_R=\frac{q}{d_B},
\qquad
\tau_1=d_Bt_1,
\qquad
\tau_2=d_Bt_2.
\]

Therefore
\[
\boxed{
\mu_R(\tau)
=
\mu_0
+
A_R
\left[
1+
\cos\left(
\frac{2\pi \tau}{T_R}-\phi
\right)
\right]
+
\frac{D_R}{2}
\left[
\tanh(q_R(\tau-\tau_1))
-
\tanh(q_R(\tau-\tau_2))
\right].
}
\]

For simplicity, we usually write \(q\) instead of \(q_R\) after nondimensionalization.

\subsection{7. CYP26A1 equation}

The dimensional equation is
\[
\frac{dC}{dt}
=
a_C
+
v_{RC}\frac{R}{K_{RC}+R}
+
v_{BC}
\frac{B^{n_B}}
{K_{BC}^{n_B}+B^{n_B}}
-
d_CC.
\]

Substitute
\[
C=C_0c,\qquad R=R_0r,\qquad B=B_0b.
\]

Then
\[
d_BC_0\frac{dc}{d\tau}
=
a_C
+
v_{RC}\frac{R_0r}{K_{RC}+R_0r}
+
v_{BC}
\frac{B_0^{n_B}b^{n_B}}
{K_{BC}^{n_B}+B_0^{n_B}b^{n_B}}
-
d_CC_0c.
\]

Divide by \(d_BC_0\):
\[
\frac{dc}{d\tau}
=
\frac{a_C}{d_BC_0}
+
\frac{v_{RC}}{d_BC_0}
\frac{r}{\frac{K_{RC}}{R_0}+r}
+
\frac{v_{BC}}{d_BC_0}
\frac{b^{n_B}}
{\left(\frac{K_{BC}}{B_0}\right)^{n_B}+b^{n_B}}
-
\frac{d_C}{d_B}c.
\]

Define
\[
\varepsilon_C=\frac{d_B}{d_C},
\qquad
a_C^*=\frac{a_C}{d_CC_0},
\]
\[
\eta_{RC}=\frac{v_{RC}}{d_CC_0},
\qquad
\kappa_{RC}=\frac{K_{RC}}{R_0},
\]
\[
\eta_{BC}=\frac{v_{BC}}{d_CC_0},
\qquad
\kappa_{BC}=\frac{K_{BC}}{B_0}.
\]

The nondimensional equation is
\[
\boxed{
\frac{dc}{d\tau}
=
\frac{1}{\varepsilon_C}
\left[
a_C^*
+
\eta_{RC}\frac{r}{\kappa_{RC}+r}
+
\eta_{BC}
\frac{b^{n_B}}
{\kappa_{BC}^{n_B}+b^{n_B}}
-
c
\right].
}
\]

If \(C_0=a_C/d_C\), then \(a_C^*=1\). In simulations we keep \(a_C^*\) as a free dimensionless basal CYP production parameter.

\section{Final Nondimensional System}

The complete nondimensional model is
\[
\boxed{
\frac{db}{d\tau}
=
w
+
\eta_{13}
\frac{h_{13}^{n_H}}
{\kappa_{13}^{n_H}+h_{13}^{n_H}}
-
b
-
\lambda_Ppb
-
\lambda_5\frac{h_5b}{\kappa_5+b}
}
\]

\[
\boxed{
\frac{dp}{d\tau}
=
\frac{1}{\varepsilon_P}
\left[
\frac{1+\rho_5h_5}
{1+\rho_Bb+\rho_{13}h_{13}}
-
\delta_P(\theta_P)p
\right]
}
\]

\[
\boxed{
\frac{dh_5}{d\tau}
=
\frac{1}{\varepsilon_5}
\left[
a_5^*
+
\eta_R\frac{r}{\kappa_R+r}
-
h_5
-
\eta_M\frac{mh_5}{\kappa_M+m}
\right]
}
\]

\[
\boxed{
\frac{dh_{13}}{d\tau}
=
\frac{1}{\varepsilon_{13}}
\left[
a_{13}^*
+
\eta_{B13}
\frac{b^{n_B}}
{\kappa_{B13}^{n_B}+b^{n_B}}
+
\eta_{M13}
\frac{m^{n_M}}
{\kappa_{M13}^{n_M}+m^{n_M}}
-
h_{13}
\right]
}
\]

\[
\boxed{
\frac{dm}{d\tau}
=
\frac{1}{\varepsilon_M}
\left[
a_M^*
+
\eta_{BM}
\frac{b^{n_B}}
{\kappa_{BM}^{n_B}+b^{n_B}}
-
m
\right]
}
\]

\[
\boxed{
\frac{dr}{d\tau}
=
\frac{1}{\varepsilon_R}
\left[
\mu_R(\tau)
-
r
-
\lambda_Ccr
\right]
}
\]

\[
\boxed{
\frac{dc}{d\tau}
=
\frac{1}{\varepsilon_C}
\left[
a_C^*
+
\eta_{RC}\frac{r}{\kappa_{RC}+r}
+
\eta_{BC}
\frac{b^{n_B}}
{\kappa_{BC}^{n_B}+b^{n_B}}
-
c
\right].
}
\]

The APC mutation function is
\[
\boxed{
\delta_P(\theta_P)=1+\delta_{P1}(1-\theta_P).
}
\]

Here,
\[
\theta_P=1
\]
corresponds to normal APC function, while
\[
0<\theta_P<1
\]
corresponds to partial APC loss or APC-mutant behavior.

The nondimensional RA input is
\[
\boxed{
\mu_R(\tau)
=
\mu_0
+
A_R
\left[
1+
\cos\left(
\frac{2\pi \tau}{T_R}-\phi
\right)
\right]
+
\frac{D_R}{2}
\left[
\tanh(q(\tau-\tau_1))
-
\tanh(q(\tau-\tau_2))
\right].
}
\]

\section{Stemness Measures}

The dynamic stemness index is defined by
\[
\boxed{
S(\tau)
=
\frac{
b(\tau)\left(1+\alpha_{13}h_{13}(\tau)\right)
}
{
\left(1+p(\tau)\right)
\left(1+\alpha_5h_5(\tau)\right)
}.
}
\]

The equilibrium stemness index is
\[
\boxed{
S^*
=
\frac{
b^*\left(1+\alpha_{13}h_{13}^*\right)
}
{
\left(1+p^*\right)
\left(1+\alpha_5h_5^*\right)
},
}
\]
where
\[
(b^*,p^*,h_5^*,h_{13}^*,m^*,r^*,c^*)
\]
is the steady-state solution of the nondimensional system.
 Next step is to do the following:
  Solve the forward problem, sensitivity, bifurcation, implement PINNs, stochastic etc.


\newpage

\subsection{Dimensional Interpretation of the Dimensionless Model}

The model was formulated and simulated in dimensionless form in order to compare
the relative strengths of the regulatory interactions among WNT, HOX, RA, and
CYP signaling. To provide a biologically interpretable dimensional realization of
the simulations, we introduce the reference scales
\[
d_B=1\ \mathrm{hr}^{-1},
\]
and
\[
B_0=0.20,\quad
P_0=1.00,\quad
H_{5,0}=0.80,\quad
H_{13,0}=0.30,\quad
M_0=0.30,\quad
R_0=0.60,\quad
C_0=0.40.
\]
These values correspond to the initial conditions used in the numerical simulations:
\[
y_0=(0.20,\ 1.00,\ 0.80,\ 0.30,\ 0.30,\ 0.60,\ 0.40),
\]
where the state variables are ordered as
\[
(B,\ P,\ H_5,\ H_{13},\ M,\ R,\ C).
\]

With the reference degradation rate
\[
d_B=1~\mathrm{hr}^{-1},
\]
the dimensionless time variable is defined by
\[
\tau=d_Bt.
\]
Consequently, when the dimensional time \(t\) is expressed in hours, the numerical
value of \(\tau\) is equal to the numerical value of \(t\). For example,
\(\tau=40\) corresponds to \(t=40~\mathrm{hr}\), and \(\tau=88\) corresponds to
\(t=88~\mathrm{hr}\). Thus, one dimensionless time unit represents one hour on
the adopted biological time scale.

The dimensional parameter values reported below are estimated values based on
the selected reference scales and the dimensionless parameter set. They are
chosen such that nondimensionalization reproduces exactly the parameter values
used in the numerical simulations. Consequently, the dimensional and
dimensionless formulations are mathematically equivalent, while the dimensional
representation provides a biologically interpretable description of the model
parameters and time scales.

This interpretation is appropriate for an in vitro signaling model describing the
regulation of the WNT, HOX, and retinoic acid pathways in colorectal cancer.
To distinguish treatment-induced responses from the initial transient associated
with the prescribed initial conditions, the model is first allowed to evolve under
untreated conditions for 40 hours before ATRA administration. The treatment is
then applied over the interval
\[
40 \leq t \leq 88~\mathrm{hr},
\]
which is equivalent to
\[
40 \leq \tau \leq 88.
\]
This corresponds to a 48-hour ATRA exposure following a 40-hour equilibration
period. Such treatment durations are consistent with widely used in vitro protocols
for investigating the effects of all-trans retinoic acid on colorectal cancer cell
proliferation, differentiation, stemness, and signaling
pathways~\cite{ordonezmoran2015hoxa5,facey2020retinoids,xu2022atra}. Accordingly, the treatment
window should be interpreted as a controlled cell-culture exposure designed to
study pathway dynamics rather than as a clinical treatment schedule.
\begin{table}[ht]
\centering
\caption{Reference scales used to recover the dimensional parameter set from the
dimensionless model. These scales preserve exactly the dimensionless parameter
values employed in all numerical simulations.}
\label{tab:reference_scales}
\begin{tabular}{lll}
\toprule
Reference quantity & Description & Value \\
\midrule
$d_B$ & Reference degradation rate of $\beta$-catenin & $1.00~\mathrm{hr}^{-1}$\\
$B_0$ & Reference $\beta$-catenin concentration & $0.20$\\
$P_0$ & Reference APC concentration & $1.00$\\
$H_{5,0}$ & Reference HOXA5 concentration & $0.80$\\
$H_{13,0}$ & Reference HOXA13 concentration & $0.30$\\
$M_0$ & Reference MYC concentration & $0.30$\\
$R_0$ & Reference RA:RAR concentration & $0.60$\\
$C_0$ & Reference CYP26A1 concentration & $0.40$\\
\bottomrule
\end{tabular}
\end{table}

\begin{longtable}{lll}
\caption{Recovered dimensional parameter values. These values constitute one
consistent dimensional realization of the dimensionless model and reproduce
exactly the numerical simulations after nondimensionalization.}
\label{tab:dimensional_parameters}
\\

\toprule
Parameter & Biological interpretation / Units & Value\\
\midrule
\endfirsthead

\caption[]{Recovered dimensional parameter values (continued).}\\
\toprule
Parameter & Biological interpretation / Units & Value\\
\midrule
\endhead

\bottomrule
\endfoot

$d_B$ & $\beta$-catenin degradation rate $(\mathrm{hr}^{-1})$ & 1.000\\
$v_{13}$ & HOXA13 synthesis rate $(\mathrm{hr}^{-1})$ & 0.150\\
$K_{13}$ & HOXA13 activation constant & 0.165\\
$k_{PB}$ & APC inhibition coefficient $(\mathrm{hr}^{-1})$ & 1.600\\
$v_5$ & HOXA5 synthesis rate $(\mathrm{hr}^{-1})$ & 0.325\\
$K_5$ & $\beta$-catenin activation constant & 0.100\\
$\gamma_{P0}$ & Basal APC degradation rate $(\mathrm{hr}^{-1})$ & 1.000\\
$\gamma_{P1}$ & Enhanced APC degradation rate $(\mathrm{hr}^{-1})$ & 3.500\\
$\rho_5$ & HOXA5 repression coefficient & 1.375\\
$\rho_B$ & $\beta$-catenin repression coefficient & 5.500\\
$\rho_{13}$ & HOXA13 activation coefficient & 4.333\\
$d_5$ & HOXA5 degradation rate $(\mathrm{hr}^{-1})$ & 0.833\\
$a_5$ & Basal HOXA5 production $(\mathrm{hr}^{-1})$ & 0.100\\
$v_R$ & RA-induced HOXA5 synthesis $(\mathrm{hr}^{-1})$ & 1.667\\
$K_R$ & RA activation constant & 0.240\\
$v_M$ & MYC-induced HOXA5 synthesis $(\mathrm{hr}^{-1})$ & 1.667\\
$K_M$ & MYC activation constant & 0.150\\
$d_{13}$ & HOXA13 degradation rate $(\mathrm{hr}^{-1})$ & 1.000\\
$a_{13}$ & Basal HOXA13 synthesis $(\mathrm{hr}^{-1})$ & 0.054\\
$v_{B13}$ & $\beta$-catenin-induced HOXA13 synthesis $(\mathrm{hr}^{-1})$ & 0.285\\
$K_{B13}$ & $\beta$-catenin activation constant & 0.100\\
$v_{M13}$ & MYC-induced HOXA13 synthesis $(\mathrm{hr}^{-1})$ & 0.165\\
$K_{M13}$ & MYC activation constant & 0.150\\
$d_M$ & MYC degradation rate $(\mathrm{hr}^{-1})$ & 1.667\\
$a_M$ & Basal MYC synthesis $(\mathrm{hr}^{-1})$ & 0.090\\
$v_{BM}$ & $\beta$-catenin-induced MYC synthesis $(\mathrm{hr}^{-1})$ & 0.675\\
$K_{BM}$ & $\beta$-catenin activation constant & 0.100\\
$d_R$ & RA degradation rate $(\mathrm{hr}^{-1})$ & 2.500\\
$u_0$ & Basal RA production $(\mathrm{hr}^{-1})$ & 0.525\\
$A_d$ & Dietary RA input $(\mathrm{hr}^{-1})$ & 0.060\\
$D$ & ATRA treatment amplitude $(\mathrm{hr}^{-1})$ & 2.250\\
$T_d$ & Dietary RA period (hr) & 24\\
$q$ & Treatment transition rate $(\mathrm{hr}^{-1})$ & 0.300\\
$t_1$ & Treatment start time (hr) & 40\\
$t_2$ & Treatment end time (hr) & 80\\
$k_{CR}$ & CYP-mediated RA degradation $(\mathrm{hr}^{-1})$ & 5.313\\
$d_C$ & CYP degradation rate $(\mathrm{hr}^{-1})$ & 1.250\\
$a_C$ & Basal CYP synthesis $(\mathrm{hr}^{-1})$ & 0.040\\
$v_{RC}$ & RA-induced CYP synthesis $(\mathrm{hr}^{-1})$ & 0.750\\
$K_{RC}$ & RA activation constant & 0.300\\
$v_{BC}$ & $\beta$-catenin-induced CYP synthesis $(\mathrm{hr}^{-1})$ & 0.750\\
$K_{BC}$ & $\beta$-catenin activation constant & 0.100\\

\end{longtable}

\begin{longtable}{lll}
\caption{Dimensionless parameter values used throughout the numerical simulations.}
\label{tab:dimensionless_parameters}
\\

\toprule
Parameter & Definition & Value\\
\midrule
\endfirsthead

\caption[]{Dimensionless parameters (continued).}\\
\toprule
Parameter & Definition & Value\\
\midrule
\endhead

\bottomrule
\endfoot

$\eta_{13}$ & $\dfrac{v_{13}}{d_BB_0}$ & 0.75\\
$\kappa_{13}$ & $\dfrac{K_{13}}{H_{13,0}}$ & 0.55\\
$\lambda_P$ & $\dfrac{k_{PB}P_0}{d_B}$ & 1.60\\
$\lambda_5$ & $\dfrac{v_5H_{5,0}}{d_BB_0}$ & 1.30\\
$\kappa_5$ & $\dfrac{K_5}{B_0}$ & 0.50\\
$\epsilon_P$ & $\dfrac{d_B}{\gamma_{P0}}$ & 1.00\\
$\delta_{P1}$ & $\dfrac{\gamma_{P1}}{\gamma_{P0}}$ & 3.50\\
$\epsilon_5$ & $\dfrac{d_B}{d_5}$ & 1.20\\
$\eta_R$ & $\dfrac{v_R}{d_5H_{5,0}}$ & 2.50\\
$\kappa_R$ & $\dfrac{K_R}{R_0}$ & 0.40\\
$\eta_M$ & $\dfrac{v_M}{d_5H_{5,0}}$ & 2.50\\
$\kappa_M$ & $\dfrac{K_M}{M_0}$ & 0.50\\
$\epsilon_{13}$ & $\dfrac{d_B}{d_{13}}$ & 1.00\\
$\eta_{B13}$ & $\dfrac{v_{B13}}{d_{13}H_{13,0}}$ & 0.95\\
$\kappa_{B13}$ & $\dfrac{K_{B13}}{B_0}$ & 0.50\\
$\eta_{M13}$ & $\dfrac{v_{M13}}{d_{13}H_{13,0}}$ & 0.55\\
$\kappa_{M13}$ & $\dfrac{K_{M13}}{M_0}$ & 0.50\\
$\epsilon_M$ & $\dfrac{d_B}{d_M}$ & 0.60\\
$\eta_{BM}$ & $\dfrac{v_{BM}}{d_MM_0}$ & 1.35\\
$\kappa_{BM}$ & $\dfrac{K_{BM}}{B_0}$ & 0.50\\
$\epsilon_R$ & $\dfrac{d_B}{d_R}$ & 0.40\\
$\mu_0$ & $\dfrac{u_0}{d_RR_0}$ & 0.35\\
$A_R$ & $\dfrac{A_d}{d_RR_0}$ & 0.04\\
$D_R$ & $\dfrac{D}{d_RR_0}$ & 1.50\\
$T_R$ & $d_BT_d$ & 24\\
$q_\tau$ & $\dfrac{q}{d_B}$ & 0.30\\
$\tau_1,\tau_2$ & $d_Bt_1,\ d_Bt_2$ & 40, 80\\
$\lambda_C$ & $\dfrac{k_{CR}C_0}{d_R}$ & 0.85\\
$\epsilon_C$ & $\dfrac{d_B}{d_C}$ & 0.80\\
$\eta_{RC}$ & $\dfrac{v_{RC}}{d_CC_0}$ & 1.50\\
$\kappa_{RC}$ & $\dfrac{K_{RC}}{R_0}$ & 0.50\\
$\eta_{BC}$ & $\dfrac{v_{BC}}{d_CC_0}$ & 1.50\\
$\kappa_{BC}$ & $\dfrac{K_{BC}}{B_0}$ & 0.50\\

$W$ & WNT stimulus & $0.80,\ 1.00,\ 1.50,\ 2.00$\\
$\Theta_P$ & APC reduction factor & $1.00,\ 0.75,\ 0.50,\ 0.25$\\

$n_B$ & Hill coefficient & 2\\
$n_M$ & Hill coefficient & 2\\
$n_H$ & Hill coefficient & 2\\
$\alpha_5$ & Stemness weight & 1\\
$\alpha_{13}$ & Stemness weight & 1\\

\end{longtable}


\subsection{Biological Interpretation of the Time Scale}

The choice
\[
d_B=1\ \mathrm{hr}^{-1}
\]
places the model on an hourly signaling scale. This is reasonable for a reduced
ODE model of pathway-level regulation in cultured colorectal cancer cells, because
the variables represent effective pathway activities rather than direct measurements
of tumor growth over clinical time. Under this scaling, the periodic RA forcing
\[
T_R=24
\]
corresponds to a 24-hour cycle, while the ATRA treatment window
\[
40\leq \tau \leq 80
\]
corresponds to a 40-hour exposure.

This time scale is consistent with experimental studies in which ATRA and related
retinoids are used to probe colorectal cancer cell differentiation, stemness, and
signaling responses over short in vitro windows~\cite{ordonezmoran2015hoxa5,facey2020retinoids,xu2022atra}.
In particular, Ordóñez-Morán et al. showed that HOXA5 counteracts intestinal
stem-cell traits by inhibiting WNT signaling in colorectal cancer, and their
supplementary results indicate that retinoids such as ATRA can induce HOXA5
expression~\cite{ordonezmoran2015hoxa5}. Facey and Boman reviewed evidence that RA
signaling is altered in colorectal cancer and that retinoids regulate colon stem-cell
differentiation and colorectal cancer stemness~\cite{facey2020retinoids}. Xu et al.
further showed that ATRA suppresses colon cancer cell activity through regulation
of the sonic hedgehog pathway~\cite{xu2022atra}. These studies support the use of
a short-term ATRA exposure window in the present model.

Therefore, the simulated ATRA window should be interpreted as a cell-culture
treatment exposure designed to study pathway response, not as a direct clinical
treatment duration. The dimensionalization preserves the simulated dimensionless
model while providing a biologically interpretable hourly time scale.


\newpage
\section{Scientific Machine Learning Extensions}

The deterministic WNT--HOX--RA model provides a mechanistic framework for
studying colorectal cancer signaling, stemness, and treatment response. The model
has the general form
\[
\frac{dy}{d\tau}=f(y,\theta),
\]
where
\[
y=(B,P,H_5,H_{13},M,R,C)
\]
denotes the state vector of $\beta$-catenin, APC, HOXA5, HOXA13, MYC, RA, and
CYP26A1, and $\theta$ denotes the vector of model parameters.

The following three projects extend the model using scientific machine learning.
Each project preserves the biological structure of the mechanistic model while
addressing a different limitation: parameter uncertainty, missing regulatory
mechanisms, and biological stochasticity.

\subsection{Project 1: Bayesian Physics-Informed Neural Networks for Parameter Uncertainty}

The current inverse PINN framework estimates a single optimal value for each
unknown parameter. However, biological parameters are uncertain because they are
inferred from noisy and limited experimental observations. The goal of this project
is to estimate probability distributions for selected model parameters rather than
single point estimates.

The Bayesian inverse problem is formulated as
\[
p(\theta\mid D)\propto p(D\mid \theta)p(\theta),
\]
where $D$ denotes the observed data, $p(D\mid \theta)$ is the likelihood, and
$p(\theta)$ is the prior distribution.

A suitable set of parameters to estimate is
\[
\Theta=
\{\eta_R,\eta_{BM},\eta_{B13},\lambda_C,\theta_P\}.
\]
These parameters are biologically meaningful because they regulate RA-induced
HOXA5 activation, $\beta$-catenin-induced MYC activation, $\beta$-catenin-induced
HOXA13 activation, CYP-mediated RA degradation, and APC mutation severity.

The neural network takes time as input and predicts all seven state variables:
\[
\tau \longmapsto
(B(\tau),P(\tau),H_5(\tau),H_{13}(\tau),M(\tau),R(\tau),C(\tau)).
\]
A fully connected neural network with smooth activation functions such as
$\tanh$ or Swish/SiLU may be used.

The loss function has the form
\[
\mathcal L
=
\mathcal L_{\mathrm{data}}
+
\mathcal L_{\mathrm{physics}}
+
\mathcal L_{\mathrm{prior}}.
\]
Here, $\mathcal L_{\mathrm{data}}$ measures agreement with observations,
$\mathcal L_{\mathrm{physics}}$ enforces the ODE residuals, and
$\mathcal L_{\mathrm{prior}}$ incorporates prior biological information on the
parameters.

The main outputs of this project are posterior distributions, credible intervals,
and uncertainty bands for model predictions. This allows one to answer the
biological question: which parameters can be estimated confidently from available
data, and which remain uncertain?

\subsection{Project 2: Hybrid Mechanistic--Neural Model for Learning Missing Regulation}

Mechanistic models encode known biological interactions, but some regulatory
relationships are often incomplete or uncertain. The goal of this project is to
preserve the known WNT--HOX--RA mechanisms while using a neural network to learn
a missing regulatory contribution.

The hybrid model is written as
\[
\frac{dy}{d\tau}
=
f_{\mathrm{known}}(y,\theta)
+
f_{\mathrm{NN}}(y).
\]
The term $f_{\mathrm{known}}$ contains the mechanistic ODE model, while
$f_{\mathrm{NN}}$ represents a small neural correction term.

For this model, a biologically meaningful choice is to introduce the neural
correction into the MYC equation:
\[
\frac{dM}{d\tau}
=
\frac{1}{\epsilon_M}
\left[
a_M+\eta_{BM}\frac{B^{n_B}}{\kappa_{BM}^{n_B}+B^{n_B}}-M
\right]
+
f_{\mathrm{NN}}(H_{13},R).
\]
Here, $f_{\mathrm{NN}}(H_{13},R)$ represents an unknown regulatory contribution
from HOXA13 and RA to MYC dynamics. This is a useful candidate because HOXA13 and
RA both influence proliferation and differentiation, but their combined indirect
effect on MYC is not fully characterized.

The neural network has the structure
\[
(H_{13},R)\longmapsto f_{\mathrm{NN}}(H_{13},R).
\]
A small network is preferred, for example four hidden layers with 32--64 neurons
per layer and a smooth activation function such as Swish/SiLU. To preserve
interpretability, the correction term should be regularized using
\[
\lambda_{\mathrm{NN}}\|f_{\mathrm{NN}}\|^2.
\]
The training loss is
\[
\mathcal L
=
\mathcal L_{\mathrm{data}}
+
\mathcal L_{\mathrm{physics}}
+
\lambda_{\mathrm{NN}}\|f_{\mathrm{NN}}\|^2.
\]

After training, the learned surface
\[
f_{\mathrm{NN}}(H_{13},R)
\]
should be plotted and interpreted biologically. This allows the model to generate
new hypotheses about how HOXA13 and RA may jointly regulate MYC. The hybrid model
should be compared against the original mechanistic model using trajectory error,
parameter recovery, cross-validation, and information criteria such as AIC or BIC.

\subsection{Project 3: Stochastic PINNs for Biological Variability}

The deterministic model predicts the average behavior of the signaling network.
However, gene regulation and treatment response are inherently variable due to
transcriptional noise, heterogeneous drug uptake, and cell-to-cell variability.
The goal of this project is to extend the model to include stochastic dynamics.

The deterministic model
\[
\frac{dy}{d\tau}=f(y,\theta)
\]
is replaced by the stochastic differential equation
\[
dy=f(y,\theta)\,d\tau+G(y)\,dW_\tau,
\]
where $W_\tau$ is a Wiener process and $G(y)$ controls the strength of the noise.

Noise should not be added to every equation initially. A biologically reasonable
starting point is to introduce stochasticity into the RA and MYC equations:
\[
dR=f_R(y,\theta)\,d\tau+\sigma_R\,dW_R,
\]
\[
dM=f_M(y,\theta)\,d\tau+\sigma_M\,dW_M.
\]
RA is a natural stochastic variable because it is influenced by dietary input,
ATRA treatment, and CYP-mediated degradation. MYC is also a reasonable stochastic
variable because transcriptional regulators often display strong cell-to-cell
variability.

The first objective is to estimate the constant noise intensities
\[
\sigma_R,\qquad \sigma_M.
\]
A more advanced extension is to allow RA noise to depend on RA and CYP levels:
\[
\sigma_R=\sigma_{\mathrm{NN}}(R,C),
\]
where
\[
(R,C)\longmapsto \sigma_{\mathrm{NN}}(R,C)
\]
is learned by a neural network. This would allow the model to test whether
CYP26A1 activity increases or suppresses variability in RA signaling.

Numerically, stochastic trajectories can be generated using the Euler--Maruyama
method:
\[
y_{n+1}
=
y_n
+
f(y_n,\theta)\Delta \tau
+
G(y_n)\Delta W_n,
\]
where
\[
\Delta W_n\sim \mathcal N(0,\Delta \tau).
\]

The main outputs are ensembles of trajectories, mean responses, variance bands,
and probabilistic stemness predictions. Instead of reporting only one stemness
curve, this project estimates
\[
\mathbb P(S>S_{\mathrm{crit}}),
\]
the probability that the stemness index exceeds a biologically significant
threshold.

\subsection{Expected Contributions}

These three projects form a coherent scientific machine learning program for the
WNT--HOX--RA colorectal cancer model. Project 1 quantifies uncertainty in key
biological parameters using Bayesian PINNs. Project 2 learns missing regulatory
structure while preserving the mechanistic model. Project 3 introduces biological
variability and enables probabilistic prediction of treatment response.

Together, these extensions move the model from deterministic simulation toward an
uncertainty-aware, data-driven, and biologically interpretable framework for
studying colorectal cancer signaling and treatment response.


\begin{align*}
\frac{dB}{dt}
&=
a_BW
+
v_{13}\frac{H_{13}^{n_H}}{K_{13}^{n_H}+H_{13}^{n_H}}
-
d_BB
-
k_{PB}PB
-
v_5\frac{H_5B}{K_5+B},
\\[8pt]
\frac{dP}{dt}
&=
a_P
\frac{1+\rho_5H_5}
{1+\rho_BB+\rho_{13}H_{13}}
-
\Big(d_{P0}+d_{P1}(1-\Theta_P)\Big)P,
\\[8pt]
\frac{dH_5}{dt}
&=
a_5
+
v_R\frac{R}{K_R+R}
-
d_5H_5
-
v_M\frac{MH_5}{K_M+M},
\\[8pt]
\frac{dH_{13}}{dt}
&=
a_{13}
+
v_{B13}\frac{B^{n_B}}{K_{B13}^{n_B}+B^{n_B}}
+
v_{M13}\frac{M^{n_M}}{K_{M13}^{n_M}+M^{n_M}}
-
d_{13}H_{13},
\\[8pt]
\frac{dM}{dt}
&=
a_M
+
v_{BM}\frac{B^{n_B}}{K_{BM}^{n_B}+B^{n_B}}
-
d_MM,
\\[8pt]
\frac{dR}{dt}
&=
u_0
+
A_d
\left[
1+
\cos\left(
\frac{2\pi t}{T_d}-\phi
\right)
\right]
+
\frac{D}{2}
\bigg[
\tanh(q(t-t_1))
-
\tanh(q(t-t_2))
\bigg]
-
d_RR
-
k_{CR}CR,
\\[8pt]
\frac{dC}{dt}
&=
a_C
+
v_{RC}\frac{R}{K_{RC}+R}
+
v_{BC}\frac{B^{n_B}}{K_{BC}^{n_B}+B^{n_B}}
-
d_CC.
\end{align*}

\[
u_R(t)
=
u_0
+
\frac{D}{2}
\left[
\tanh(q(t-t_1))
-
\tanh(q(t-t_2))
\right].
\]


\begin{align*}
\frac{db}{d\tau}
&=
w
+
\eta_{13}\frac{h_{13}^{n_H}}{\kappa_{13}^{n_H}+h_{13}^{n_H}}
-
b
-
\lambda_Ppb
-
\lambda_5\frac{h_5b}{\kappa_5+b},
\\[8pt]
\frac{dp}{d\tau}
&=
\frac{1}{\varepsilon_P}
\left[
\frac{1+\rho_5h_5}
{1+\rho_Bb+\rho_{13}h_{13}}
-
\delta_P(\theta_P)p
\right],
\\[8pt]
\frac{dh_5}{d\tau}
&=
\frac{1}{\varepsilon_5}
\left[
a_5
+
\eta_R\frac{r}{\kappa_R+r}
-
h_5
-
\eta_M\frac{mh_5}{\kappa_M+m}
\right],
\\[8pt]
\frac{dh_{13}}{d\tau}
&=
\frac{1}{\varepsilon_{13}}
\left[
a_{13}
+
\eta_{B13}\frac{b^{n_B}}{\kappa_{B13}^{n_B}+b^{n_B}}
+
\eta_{M13}\frac{m^{n_M}}{\kappa_{M13}^{n_M}+m^{n_M}}
-
h_{13}
\right],
\\[8pt]
\frac{dm}{d\tau}
&=
\frac{1}{\varepsilon_M}
\left[
a_M
+
\eta_{BM}\frac{b^{n_B}}{\kappa_{BM}^{n_B}+b^{n_B}}
-
m
\right],
\\[8pt]
\frac{dr}{d\tau}
&=
\frac{1}{\varepsilon_R}
\left[
\mu_R(\tau)
-
r
-
\lambda_Ccr
\right],
\\[8pt]
\frac{dc}{d\tau}
&=
\frac{1}{\varepsilon_C}
\left[
a_C
+
\eta_{RC}\frac{r}{\kappa_{RC}+r}
+
\eta_{BC}\frac{b^{n_B}}{\kappa_{BC}^{n_B}+b^{n_B}}
-
c
\right].
\end{align*}

\[
\delta_P(\theta_P)=1+\delta_{P1}(1-\theta_P).
\]

\[
A_d \left[ 1 + \cos\left( \frac{2\pi t}{T_d} - \phi \right) \right]
\]


\newpage
\begin{abstract}
Reduced form\dots{} (250 words max)
\end{abstract}

intro - summarize papers we have read; biology + PINN; then Amir, Nate, PK use A, B, C, D to do the techniques from the papers, discussions, etc.

Then, modeling, schematic, assumptions \\
Then equations, non-dimensionalization

Numerical results; sci-py - talk about meaning of spikes or dips biologically

PINN results; algorithm, why did we use PINN to approximate solution; advantages
- errors, convergence, etc.

also sensitivity, identifiability, bayesian PINN, etc.

use first presentation for reference

Last part - Discussion - restate what we talked about, discuss importance of our model, limitations of the model - and future direction (Hybrid Mechanistic, Stochastic Modeling)\dots{} must incorporate biological significance to these future methods.
Modeling at gene level is stochastic in nature.
Hybrid mechanistic - we don't know all interactions between things; modeling simplifies interactions and may miss things; learn behaviors from data. But we need real data instead of synthetic data; hard to get data from normal colon~\cite{liu2022hybrid,yazdani2020systems}


\bibliographystyle{unsrt}
\bibliography{references}

\end{document}

